\documentclass[a4paper]{scrartcl}
\usepackage{amsmath,amssymb,amsthm}
\usepackage{bm}
\usepackage{biblatex}
\usepackage{caption}
\usepackage{float}
\usepackage[colorlinks=true, allcolors=black, urlcolor=cyan]{hyperref}
\usepackage{graphicx}
\usepackage[framed,numbered]{matlab-prettifier}
\usepackage{mathtools}
\usepackage{mcode}
\usepackage{physics}
\usepackage{siunitx}
\usepackage{cancel}

\renewcommand{\thesection}{}
\renewcommand{\thesubsection}{\arabic{section}.\alph{subsection}}

\title{Assignment 2}
\subtitle{ELEC 442 - Introduction to Robotics}
\author{Sondre Myrberg (81113433) \and Ola Helbaek (68776772)}


\setlength{\parindent}{0pt} % Disable indentation
\mathtoolsset{showonlyrefs} % Only show equation numbers on referenced equations

\def\undertilde#1{\mathord{\vtop{\ialign{##\crcr
$\hfil\displaystyle{#1}\hfil$\crcr\noalign{\kern1.5pt\nointerlineskip}
$\hfil\widetilde{}\hfil$\crcr\noalign{\kern1.5pt}}}}} %Lousy way, but need this undertilde



\begin{document}

\hypersetup{pageanchor=false}
\begin{titlepage}
    \maketitle
    \vfill
    \vfill
    \vfill
    \vfill
    \includegraphics[width=0.95\textwidth]{../../ubc_logo.pdf}
    \vfill
    \vfill
\end{titlepage}
\hypersetup{pageanchor=true}

\section{Solving the inverse kinematics problem for the PUMA 560 robot}

For the first angle, $\theta_1$ we use that the centre of the spherical wrist, $\undertilde{o_4}$, compared to the $xy$-plane is only dependent on $\theta_1$. We also know that for $\theta_1 = 0$ $\undertilde{o_4}$ must be in the plane defined by $y - 149.09 \si{\mm} = 0$. This gives us the relation between $u$ and $w$, where $w$ is the vector between the origin and $\undertilde{o_4}$ and $u$ is the vector to where $\undertilde{o_4}$ would be if $\theta_1 = 0$, such that we kan use the inplemented function \mcode{KahanP2}. The relation is given as
\begin{equation}
    \begin{aligned}
        \norm{w}^2 = o_{4,x}^2 + o_{4,y}^2 + o_{4,z}^2 &= u_x^2 + u_y^2 + u_z^2 = \norm{u}^2\\
        \implies o_{4,x}^2 + o_{4,y}^2 + \cancel{o_{4,z}^2} &= u_x^2 + 149.09^2 + \cancel{o_{4,z}^2} \quad &&\text{as desired $z$-value doesn't} \\
        \implies u_x &= \sqrt{o_{4,x}^2 + o_{4,y}^2 - 149.09^2}                                             &&\text{change with $\theta_1$}
    \end{aligned}
\end{equation}

Given this and that $\undertilde{o_0} = \begin{pmatrix} 0 & 0 & 0 \end{pmatrix}^\top$ we can use the implemented function \mcode{KahanP2} with $s = \bm{k}_0$, $\hat{\bm{u}} = \begin{bmatrix} u_x & 149.09 & o_{4,z} \end{bmatrix}$ and $\bm{w} = \undertilde{o_4}-\undertilde{o_0}$. The function is implemented as shown in \autoref{lst:KahanP2}. In this implementation the function itself normalizes the vectors such that we do not have to think about feeding this into the function. 

\lstinputlisting[style=Matlab-editor, caption=MATLAB implementation of the Kahan P2 problem, captionpos = b, label=lst:KahanP2]{../../MATLAB/Assignment2/KahanP2.m}

\end{document}